\section{DeepSeek、开源与闭源}

前面 Scaling Law 讨论的是能力增长，本章转向能力如何被行业感知和扩散。DeepSeek 把开源、效率和真实效果推到了行业叙事中心。

2025 年章节标题是“向 DeepSeek 学习”。张鹏提到，DeepSeek 某种程度上也教会行业：技术上不需要总靠大发布会，最终还是回归实际应用效果。DeepSeek 同时改变了开源叙事和效率叙事，让模型公司重新思考开放、生态和商业控制之间的关系。

\begin{figure}[H]
\centering
\includegraphics[width=0.94\textwidth]{open-vs-closed.png}
\caption{开源 vs 闭源：生态影响和商业控制之间的取舍。}
\end{figure}

\begin{knowledgebox}{读图：开源和闭源不是道德二分}
开源能扩大生态、获得开发者反馈、建立技术声誉；闭源能保留商业控制、服务企业客户、处理安全和合规边界。模型公司通常需要在不同产品层和不同客户层做组合。
\end{knowledgebox}

\subsection{开源带来的全球反馈}

张鹏提到智谱模型开源后，海外评价很高，甚至出现其他产品或公司直接使用、蒸馏、裁剪模型的现象。这既说明开源有影响力，也说明开源会带来被复用、被改造、被重新包装的复杂生态。

\begin{warningbox}{开源不是只收获赞美}
开源会带来开发者声誉和生态扩散，也会带来商业归属模糊、模型被二次包装、收益难捕捉等问题。开源策略必须和商业化策略一起设计。
\end{warningbox}

\subsection{本章小结}

DeepSeek 之后，模型公司的竞争更加重视效率、真实效果和开源生态。开源与闭源不是二选一，而是模型公司如何分配生态影响、商业价值和安全边界。

